\begin{enumerate}

\item La suma y diferencia de remolinos dan la entrada.

\[\fitmath{\mathrm{ct}_{2}=\mathrm{uu}\cdot \left(1-\mathrm{rea}+\frac{\psi}{2}\right)}\]

\[\fitmath{\mathrm{ct}_{2}=(324.825)\cdot \left(1-(0.46)+\frac{(1.42)}{2}\right)\approx 406.0312\,\mathrm{m/s}}\]

\item Se obtiene la descarga absoluta.

\[\fitmath{\mathrm{ct}_{3}=\mathrm{uu}\cdot \left(1-\mathrm{rea}-\frac{\psi}{2}\right)}\]

\[\fitmath{\mathrm{ct}_{3}=(324.825)\cdot \left(1-(0.46)-\frac{(1.42)}{2}\right)\approx -55.22025\,\mathrm{m/s}}\]

\item La referencia relativa resta U a la descarga.

\[\fitmath{\mathrm{wt}_{3}=\mathrm{uu}\cdot \left(-\mathrm{rea}-\frac{\psi}{2}\right)}\]

\[\fitmath{\mathrm{wt}_{3}=(324.825)\cdot \left(-(0.46)-\frac{(1.42)}{2}\right)\approx -380.0453\,\mathrm{m/s}}\]

\end{enumerate}
